\documentclass[
  ngerman,
  symmetric,
  cdfont=on,
  numbers=noenddot
]{tudscrreprt}

% Optional: Serif font like in your template
\usepackage[bitstream-charter]{mathdesign}
\addtokomafont{disposition}{\rmfamily}

\usepackage[T1]{fontenc}
\usepackage[utf8]{inputenc}
\usepackage[ngerman]{babel}
\usepackage[hidelinks]{hyperref}

\begin{document}

% Briefkopf / Titelseite (aus deiner Vorlage)
\pagenumbering{Roman}

%\faculty{Fakultät Informatik}
\faculty{Faculty of Computer Science}
\department{}
%\institute{Institut für Systemarchitektur}
\institute{Institute of Systems Architecture}
%\chair{Lehrstuhl Distributed and Networked Systems}
\chair{Chair of Distributed and Networked Systems}

\title{%
  Exposé: Entwicklung von Ansätzen zur Installation und dem Betrieb von Kubernetes-Clustern auf Rechnerknoten im Weltall
}

% Fix für englische Lokalisierung von Abschlussarbeiten (optional; nur nötig, wenn du english als babel nutzt)
\makeatletter
 \tud@localization@english{\bachelorthesisname}{Bachelor's Thesis}
\makeatother
\date{}
\author{Jannes Mueller}

\maketitle
\newpage

% Ab hier dein Exposé-Inhalt
\pagenumbering{arabic}

\section{Motivation}
Eine der wesentlichen Unterschiede zwischen der Entwicklung und den Betrieb von Satelliten und 
die Entwicklung und der Betrieb von Diensten im IT Umfeld ist die Möglichkeit über die gesamte Lebensdauer
hinweg nach Bedarf Anpassungen vorzunehmen. Die Trends der Entwicklungen von IT-Diensten der letzten Jahre 
zeigen die Vorteile kleine Anpassungen und Korrekturen im Code schnell deployen zu können. So können 
einfach zuvor gemachte Fehler ausgebessert und etwaige Änderungen in den Produktanforderungen auch 
nachträglich noch umgesetzt werden. Auf diese Weise kann sich der Lebenszyklus einer Anwendung deutlich 
verlängern. Dies ist sowohl vorteilhaft für den Kunden als auch für den Entwickler. 
Von solch einer Flexibilität und Effizienz kann bei Satelliten bislang nur geträumt werden. Hierbei handelt
es sich in der Regel um Maschinen, bei denen die Entwicklung von Software und Hardware zugleich erfolgt 
und die nach ihrem Deployment im All nur sehr eingeschränkt geupdatet werden können. 
Die Einführung von zusätzlicher oder erweiterter Funktionalität zu einem späteren Zeitpunkt ist somit 
nicht gegeben. Dies sorgt für eine verminderten Lebenszyklus der Maschinen, da aktualisierte 
Anforderungen nicht mehr mit bereits vorhandenen Maschinen umgesetzt werden können. 
Um diese eingeschränkte Funktionalität zu erweitern und die Wiederverwendbarkeit zu erhöhen ist es 
zwingend notwendig die Flexibilität der softwareseitigen Applikationen auf Satelliten zu erhöhen. 
Hierbei genügt ein einfacher Blick auf die Entwicklung von Software für terristische Anwendungen um 
Verbesserungen ausmachen zu können. 
Im Cloud-Computing wird heute hauptsächlich auf eine stark virtualisierte und containerisierte
Infrastruktur gesetzt, anstatt auf große monolithische Anwendungen. Dies ermöglicht sowohl eine deutlich
bessere Skalierung, als auch eine einfachere Wartung und ein schnelleres Deployment aktualisierter Versionen. 
Anstatt bei jeder Änderung eine komplett neue Version dem Kunden zur Verfügung zu stellen, muss hierbei 
lediglich der zugehörige Container neu deployed werden. Diese Flexibilität bringt viele Vorteile mit sich, 
welche ebenfalls bei der Anwendung in der Entwicklung von Satelliten Anwendung finden würden. Anstatt ein 
eigenes OS zu entwickeln, was nach gewisser Zeit nur noch begrenzt aktualisiert werden kann, könnte die 
Hauptfunktionalität als unabhängig als Container deployed werden. Auf diese Weise ist es leichter möglich 
die Funktionalität zu updaten oder sogar komplett neu zu entwickeln um nachträglich eine höhere Effizienz oder 
neue Funktionen zu implementieren. Auch eine Teilung der Zugriffszeit auf den Satelliten mit differenzierten 
Anwendungsfällen ist denkbar. Somit können sich Firmen zusammenschließen um gemeinsam die Kosten für die 
Entwicklung bereitzustellen. Ist der Satellit erst einmal im All deployed können beide Firmen durch die 
Separierung der Container unabhängig voneinander Berechnungen durchführen. 


\section{Problemanalyse}
Die Ausrichtung an terristischen Cloud-Infrastrukturen als Ansatz auf Satelliten bietet viele Vorteile, 
ist allerdings mit vielen Hindernissen verbunden. Eine einfache Übernahme des Konzeptes ist aufgrund der 
unterschiedlichen Kommunikationseigenschaften zwischen Erde und All so nicht möglich. Aufgrund der größeren
Distanz und der Begrenzung auf Funktechnologie zur Datenübertragen kommt es hierbei zu deutlich höheren 
Latenzen. Auch komplette Abbrüche der Verbindungen über einen längeren Zeitraum sind nicht untypisch. Bei 
solchen Bedingungen ist es nicht ohne weiteres möglich übliche terristische Kommunikationsprotokolle zu 
verwenden, da diese nicht für solch hohe Latenzen entwickelt wurden. Auch Verbindungsabbrüche über einen 
werden längeren Zeitraum werden aufgrund seltener Vorkommnisse auf der Erde nicht ausreichend abgefangen. 
Um dennoch diese Prinzipien im All zu implementieren, könnte eine Neuentwicklung mit Fokus auf dem Betrieb 
von Container auf Satelliten angestoßen werden. Allerdings ist eine Neuentwicklung zum Einen immer mit einem 
erheblichen Aufwand und Kosten verbunden. Zum Anderen gestaltet sich die Administration im späteren laufenden 
Betrieb schwierig, da bei einem neuen System keine Erfahrungswerte und dementsprechend wenig geschulte 
Fachkräfte bereitstehen. Daher ist es generell von Vorteil bereits existierende Lösungen, welche sich seit 
Jahren am Markt etabliert haben zu favoritisieren. Dabei kann nicht nur auf einen breite Erfahrungswerte bei
der Entwicklung dieser Lösungen vertraut werden, sonder es kann ebenfalls bereits auf geschultes Personal zur 
Administration dieser Systeme zurückgegriffen werden. Somit besteht die Herausforderung diese bereits
existierenden Lösung für den Einsatz im All auf Satelliten anzupassen.

\section{Zielstellung und Forschungsfragen}
Die allgemeine Zielsetzung ist die Entwicklung und anschließende Evaluierung von Kubernetes Cluster auf
Rechnerknoten im Weltall. Bei Kubernetes handelt es sich um ein sehr mächtiges Werkzeug zur 
Containerorchestrierung, welches durch seine Vielseitigkeit und xxxx sich im Cloud-Computing als de facto
Standard durchgesetzt hat. 

\section{Related Work}
% ...

\section{Lösungsansatz}
% ...

\section{Evaluation}
% ...

\section{Zeitplan}
% ...



\end{document}